% ============================ ANNEXES =================================
\chapter{Annexes}

Cette annexe rassemble des extraits du code réel du projet : la chaîne ETL, la
segmentation par apprentissage automatique et une vue décisionnelle SQL.

% ---------------------------------------------------------------------
\section{Chaîne ETL --- extraction, transformation, chargement (Python)}
Extraits du script \texttt{etl/etl\_agences.py}. La fonction \texttt{extract\_oracle}
lit les tables opérationnelles ; \texttt{transform} agrège les indicateurs par
agence ; \texttt{load} écrit le datamart en étoile.

\begin{lstlisting}[style=pycode,language=Python,caption={Extraction des sources depuis Oracle}]
def extract_oracle(dsn=None, user=None, password=None):
    import oracledb
    cn = oracledb.connect(
        user=user or os.environ.get("BTK_DB_USER", "system"),
        password=password or os.environ.get("BTK_DB_PWD", ""),
        dsn=dsn or os.environ.get("BTK_DB_DSN", "localhost:1521/FREEPDB1"))
    ag  = pd.read_sql("SELECT SK_AGENCE, LIBELLE_AGENCE, DISTRICT FROM AGENCE", cn)
    emp = pd.read_sql("SELECT SK_UTILISATEUR, SK_AGENCE, EST_GESTIONNAIRE "
                      "FROM B_UTILISATEURS", cn)
    cli = pd.read_sql("SELECT SK_CLIENT, SK_AGENCE FROM CLIENT_BTK", cn)
    obj = pd.read_sql(
        "SELECT SK_AGENCE, "
        "  NVL(SOUSCRIPTION_COMPTE_CHEQUES_OA,0)+NVL(SOUSCRIPTION_COMPTE_EPARGNES_OA,0)"
        "  +NVL(SOUSCRIPTION_COMPTE_COURANTS_OA,0) AS COMPTES, "
        "  NVL(PRODUCTION_CREDITS_CONSO_OA,0)+NVL(PRODUCTION_CREDITS_IMMO_OA,0)"
        "  +NVL(PRODUCTION_CREDITS_INVESTISSEMENT_OA,0) AS CREDITS, "
        "  NVL(EPARGNE_ADD_OA,0) AS EPARGNE FROM B_OBJECTIF", cn)
    poi = pd.read_sql("SELECT P.SK_UTILISATEUR, U.SK_AGENCE, P.STATUT "
                      "FROM POINTAGE P JOIN B_UTILISATEURS U "
                      "ON P.SK_UTILISATEUR = U.SK_UTILISATEUR", cn)
    cn.close()
    return ag, emp, cli, obj, poi
\end{lstlisting}

\begin{lstlisting}[style=pycode,language=Python,caption={Transformation (agrégation par agence) et chargement du datamart}]
def transform(agences, employes, clients, objectifs, pointages):
    eff = employes.groupby("SK_AGENCE").agg(
        effectif=("SK_UTILISATEUR", "count"),
        nb_gestionnaires=("EST_GESTIONNAIRE", "sum")).reset_index()
    cli = clients.groupby("SK_AGENCE").size().reset_index(name="nb_clients")
    obj = objectifs.groupby("SK_AGENCE").agg(
        total_comptes=("COMPTES", "sum"),
        production_credits=("CREDITS", "sum"),
        collecte_epargne=("EPARGNE", "sum")).reset_index()
    poi = pointages.assign(
            present=pointages["STATUT"].isin(["PRESENT", "RETARD"]).astype(int)) \
        .groupby("SK_AGENCE").agg(taux_presence=("present", "mean")).reset_index()
    dm = agences.merge(eff, on="SK_AGENCE", how="left") \
                .merge(cli, on="SK_AGENCE", how="left") \
                .merge(obj, on="SK_AGENCE", how="left") \
                .merge(poi, on="SK_AGENCE", how="left")
    return dm.rename(columns={"LIBELLE_AGENCE": "agence"})

def load(dm):
    dim_agence  = dm[["SK_AGENCE", "agence", "DISTRICT"]]
    fait_agence = dm[["SK_AGENCE"] + FEATURES]
    dim_agence.to_csv(os.path.join(ENTREPOT, "dim_agence.csv"), index=False)
    fait_agence.to_csv(os.path.join(ENTREPOT, "fait_agence.csv"), index=False)
    dm[["agence"] + FEATURES].to_csv(DATAMART, index=False)
\end{lstlisting}

% ---------------------------------------------------------------------
\section{Segmentation des agences (Python, scikit-learn)}
Extraits du script \texttt{clustering/segmentation\_agences.py} : la standardisation
des variables et la détermination du nombre optimal de clusters (méthode du coude
combinée au score de silhouette).

\begin{lstlisting}[style=pycode,language=Python,caption={Choix du nombre de clusters (coude + silhouette)}]
from sklearn.preprocessing import StandardScaler
from sklearn.cluster import KMeans
from sklearn.metrics import silhouette_score

X = StandardScaler().fit_transform(df[FEATURES])   # centrage-reduction

def choisir_k(X, kmax=8):
    ks = list(range(2, kmax + 1))
    inerties, silhouettes = [], []
    for k in ks:
        km = KMeans(n_clusters=k, n_init=10, random_state=42).fit(X)
        inerties.append(km.inertia_)                # methode du coude
        silhouettes.append(silhouette_score(X, km.labels_))
    return ks[int(np.argmax(silhouettes))]          # k au meilleur silhouette
\end{lstlisting}

\begin{lstlisting}[style=pycode,language=Python,caption={Comparaison des modèles et profilage des segments}]
modeles = {
    "K-Means":      KMeans(n_clusters=k, n_init=10, random_state=42),
    "Agglomeratif": AgglomerativeClustering(n_clusters=k),
    "GMM":          GaussianMixture(n_components=k, random_state=42),
    "DBSCAN":       DBSCAN(eps=1.2, min_samples=3),
}
for nom, modele in modeles.items():
    labels = modele.fit_predict(X)
    sil = silhouette_score(X, labels)
    db  = davies_bouldin_score(X, labels)
    ch  = calinski_harabasz_score(X, labels)
    print(nom, round(sil, 3), round(db, 3), round(ch))
# profil moyen de chaque segment
df["cluster"] = KMeans(n_clusters=k, n_init=10, random_state=42).fit_predict(X)
profils = df.groupby("cluster")[FEATURES].mean()
\end{lstlisting}

% ---------------------------------------------------------------------
\section{Vue décisionnelle (SQL)}
Exemple de vue analytique exposée à Power BI (\texttt{sql/v\_bi\_employes.sql}) :
elle joint les employés à leur agence et calcule le rôle métier.

\begin{lstlisting}[style=pycode,language=SQL,caption={Vue décisionnelle V\_BI\_EMPLOYES}]
CREATE OR REPLACE VIEW V_BI_EMPLOYES AS
SELECT
    U.SK_UTILISATEUR       AS ID_EMPLOYE,
    U.LIBELLE_UTILISATEUR  AS NOM_COMPLET,
    U.EMAIL_UTILISATEUR    AS EMAIL,
    CASE
        WHEN U.EST_GESTIONNAIRE = 1 THEN 'Gestionnaire'
        WHEN U.EST_GESTIONNAIRE = 0 AND U.SK_AGENCE = 31 THEN 'Hors Commercial'
        ELSE 'Conseiller'
    END AS ROLE_LABEL,
    CASE WHEN U.EST_SUSPENDU = 1 THEN 'Suspendu' ELSE 'Actif' END AS STATUT,
    A.LIBELLE_AGENCE AS NOM_AGENCE,
    A.DISTRICT       AS DISTRICT_AGENCE
FROM B_UTILISATEURS U
LEFT JOIN AGENCE A ON U.SK_AGENCE = A.SK_AGENCE;
\end{lstlisting}

% ---------------------------------------------------------------------
\section{Dépôt du code source}
Le code source complet du projet --- application Jakarta EE (entités, servlets,
services, pages JSP), scripts SQL d'initialisation, chaîne ETL et segmentation
Python, thème et modèle Power BI --- est versionné sur le dépôt Git associé à ce
rapport.
